\chapter{Introduction}\label{r04:ch:intro}

Reports~1 to~3 of this collection describe RISC-V Vector (RVV) support on the
\emph{new} waLBerla stack: \ps{}~2.x with our RVV backend, waLBerla's sweepgen code
generator and its v8 core. This report is about the \emph{other} stack, the one
waLBerla-wind~\cite{walberla-wind} is built on, and about how RVV code is generated
there.

\section{Why a second stack}

waLBerla-wind is an extension of waLBerla~\cite{walberla} for wind energy simulations:
lattice Boltzmann flow with turbines, wall models and a turbulence model. It uses only
waLBerla's \term{legacy core}: the C++ module \file{src/lbm_generated} and the code
generators \file{python/pystencils_walberla} and \file{python/lbmpy_walberla}. Those
generators are written against \ps{}~1.x. waLBerla-wind's continuous integration pins
\code{pystencils==1.4} and \code{lbmpy==1.4}~\cite{pystencils14,lbmpy}.

None of the RVV work of reports~1 to~3 runs on this stack. Our RVV backend is part of
\ps{}~2.x, whose internals (a typed intermediate representation, platform classes, an
intrinsic selection pass) do not exist in 1.4, and the sweepgen applications use the v8
core, which waLBerla-wind does not use. So getting RVV into waLBerla-wind meant working
on the legacy stack.

\ps{}~1.4 already ships an RVV instruction set of its own, \isa{rvv}, by Michael Kuron.
The decision taken at the start of this work was:
\begin{keyidea}[The plan in one sentence]
  Keep upstream's \isa{rvv} exactly as it is, and add our own instruction set
  \isa{rvv_ext} next to it: a copy of \isa{rvv} with optional extras (fixed vector
  length, LMUL, scalar operands, fused multiply-add, software prefetch), each switched
  on separately, so that upstream's code and every extra can be built and measured
  side by side in the same application.
\end{keyidea}

\section{What this report covers}

\begin{description}[style=nextline]
  \item[Part~I, the legacy stack] How the legacy stack differs from the new one
    (\cref{r04:ch:two-stacks}) and how its code generation turns a Python generator
    script into C++ sweeps during the CMake build (\cref{r04:ch:codegen}).
  \item[Part~II, RVV in \ps{}~1.4] How \ps{}~1.4 vectorizes a kernel
    (\cref{r04:ch:vectorizer}), what upstream's \isa{rvv} instruction set generates
    (\cref{r04:ch:upstream-rvv}), and what our \isa{rvv_ext} adds and how
    (\cref{r04:ch:rvv-ext}).
  \item[Part~III, applications and evidence] The two applications that build RVV
    variants, waLBerla's \code{UniformGridCPU} benchmark and a new waLBerla-wind
    benchmark app, the bug in waLBerla found on the way (\cref{r04:ch:apps}), and how
    the code is checked and measured (\cref{r04:ch:verification}).
\end{description}
Appendix~A lists every file this work added or changed, appendix~B the pitfalls of the
legacy build.

\section{Versions and branches}

All work is on branches named \code{rvv-support-legacy}, in this workspace and in three
of its submodules. When this report was written (26~November~2026) none of it was
committed yet; it describes the working trees of that day.
\begin{center}\small
\begin{tabularx}{\linewidth}{@{}l>{\raggedright\arraybackslash}p{4.6cm}>{\raggedright\arraybackslash}X@{}}
  \toprule
  \textbf{repository} & \textbf{branch starts at} & \textbf{holds}\\
  \midrule
  workspace & \code{main}, \code{9ff8e8b} & \file{docs/notes8/} (plan, progress log),
    \file{scripts/notes8/}, \file{scripts-legacy/}, the submodule pointers\\
  \file{submodules/pystencils} & tag \code{release/1.4}, \code{48474c1a} & the instruction
    set \isa{rvv_ext}, its headers and QEMU tests\\
  \file{submodules/walberla} & \code{rvv-support}, \code{42855054} & RISC-V CMake (from the
    2.x work), \code{UniformGridCPU} variants, one template fix\\
  \file{submodules/waLBerla-wind} & \code{main}, \code{9d74213} & the benchmark app
    \file{apps/BenchmarkApp_cpu_uniform}\\
  \bottomrule
\end{tabularx}
\end{center}
The pystencils submodule is the same repository as the 2.x fork of reports~1 to~3; the
2.x work stays on its branch \code{rvv-support} (last commit \code{78afe1ff}). Tag
\code{release/1.4} is the last 1.x release; the 2.0 rewrite was merged after it.

The target machine is the Banana~Pi BPI-F3 with the SpacemiT~K1 (eight X60 cores,
RVV~1.0, VLEN~256), as in reports~2 and~3. The K1 measurements were still running when
this report was written; \cref{r04:ch:verification} says what has been checked so far.
